\begin{proof}[Proof of \cref{obj:thm:integrated-process-certificate}]
Write \(L=\log(ed)\), \(m=n/8\), and \(K=\max\{2,\lfloor L/512\rfloor\}\). We first establish estimates for an arbitrary observed categorical law in the independent pilot and evaluation count experiment of \cref{obj:def:process-handle}. All constants introduced below are positive and depend at most on \(\epsilon\), unless called numerical.

For each pilot realization, every point \(u\) in the rectangle \(\prod_{\zeta\in\mathcal Z}[\ell_{\zeta j},u_{\zeta j}]\) is nonnegative and satisfies \(\|u-c_j^0\|_1\leq S_j\). Positivity and normalization of the Jackson convolution, together with \cref{obj:lem:bv-lipschitz}, therefore give the uniform bound
\[
\sup_{u\in\prod_{\zeta\in\mathcal Z}[\ell_{\zeta j},u_{\zeta j}]}
\|P_{j,\cdot}(u)-f_\cdot(c_j^0)\|_{BV([0,1])}
\leq c_{\mathrm{BV}}S_j.
\]
Write the polynomial on this rectangle in tensor Chebyshev form, with coefficients \(b_{j,\nu}\):
\[
P_{j,\lambda}(u)-f_\lambda(c_j^0)
 =\sum_{\nu\in\{0,\ldots,D\}^4}b_{j,\nu}(\lambda)
   \prod_{\zeta\in\mathcal Z}
   T_{\nu_\zeta}\!\left(\frac{u_\zeta-c^0_{\zeta j}}{R_{\zeta j}}\right),
\qquad D=2(K-1).
\]
Each coefficient is a normalized tensor Chebyshev integral of the uniformly bounded path, so \(\|b_{j,\nu}\|_{BV([0,1])}\leq16c_{\mathrm{BV}}S_j\). If \(A_h\) is the sum of the absolute monomial coefficients of \(T_h\), its recurrence gives \(A_0=A_1=1\), \(A_{h+1}\leq2A_h+A_{h-1}\), and hence \(A_h\leq(1+\sqrt2)^h\). Conversion to the centered monomial coefficients of \cref{obj:def:jf-estimator} now yields
\[
\sum_\alpha\|\beta_{j,\alpha}\|_{BV([0,1])}
 \leq16c_{\mathrm{BV}}S_j(D+1)^4(1+\sqrt2)^{4D}
 \leq c_{\mathrm{coeff}}S_j e^{12K}.
\]
The estimate holds for every pilot realization, including those outside \(G_j\): the pilot radii are strictly positive, and the convolution rectangle lies in \(\mathbb R_+^4\). Also \(L=1+\log d\geq1\) and \(K\leq2+L/512\). Direct expansion of this last bound yields
\[
24K+\frac{16K^2}{L}\leq114+\frac L{16}
 \leq115+\frac{\log d}{16},
\qquad
S_j^2e^{24K+16K^2/L}\leq e^{115}d^{1/16}S_j^2.
\]
These are the two coefficient bounds required by \cref{obj:def:process-handle}. % lean: CausalSmith.Stat.DiscreteBudgetvalueCurve.jacksonCoefficientBV_pilot_le
% lean: CausalSmith.Stat.DiscreteBudgetvalueCurve.jacksonDegree_exponential_budget

Here are the cell estimates used below. Put
\[
a_j=\sqrt{\frac{p_j}{mL}}+\frac1{mL},
\qquad
v_j=\sqrt{\frac{p_jL}{m}}+\frac Lm.
\]
Conditioning on the pilot counts, factorial moments of the independent evaluation counts give
\[
\mathbb E\!\left[U_h(N_{\zeta j};c^0_{\zeta j})\mid N'\right]
 =(q_{\zeta j}-c^0_{\zeta j})^h.
\]
On \(G_j\), the localization interval contains \(q_{\zeta j}\), and its radius satisfies
\(q_{\zeta j}/(mR_{\zeta j}^{\,2})\leq1/L\). The second factorial-moment identity is
\[
\mathbb E\!\left[U_h(N_{\zeta j};c^0_{\zeta j})^2\mid N'\right]
 =\sum_{r=0}^{h}{h\choose r}^{\!2}r!
   \left(\frac{q_{\zeta j}}m\right)^r
   (q_{\zeta j}-c^0_{\zeta j})^{2h-2r}.
\]
Using \(|q_{\zeta j}-c^0_{\zeta j}|\leq R_{\zeta j}\), \(q_{\zeta j}/(mR_{\zeta j}^2)\leq1/L\), and \({h\choose r}^2r!\leq h^{2r}/r!\), it gives
\[
\mathbb E\!\left[\left.
 \frac{U_h(N_{\zeta j};c^0_{\zeta j})^2}{R_{\zeta j}^{2h}}
 \right|N'\right]
 \leq\sum_{r=0}^{h}\frac{(h^2/L)^r}{r!}
 \leq e^{h^2/L}.
\]
For each factorial index \(\alpha\in\{0,\ldots,D\}^{\mathcal Z}\), independence of the four evaluation coordinates and the preceding moment bound give, on \(G_j\),
\[
\mathbb E\!\left[\left.
\prod_{\zeta\in\mathcal Z}
\frac{U_{\alpha_\zeta}(N_{\zeta j};c^0_{\zeta j})^2}
     {R_{\zeta j}^{2\alpha_\zeta}}
\right|N'\right]
\leq \exp\!\left(\frac{\sum_\zeta\alpha_\zeta^2}{L}\right)
\leq e^{4D^2/L}.
\]
Cauchy–Schwarz over the \((D+1)^4\) factorial indices and the coefficient envelope therefore give
\[
\mathbb E\!\left[\|Z_j-f_\cdot(c_j^0)\|_{BV([0,1])}^2\mid N'\right]
\leq 4\left(\sum_\alpha\|\beta_{j,\alpha}\|_{BV([0,1])}\right)^2
       \sum_{\alpha\in\{0,\ldots,D\}^{\mathcal Z}}e^{4D^2/L}
\leq 4c_{\mathrm{coeff}}^2S_j^2(D+1)^4
       e^{24K+16K^2/L}.
\]
Here \(D+1\leq2K\leq e^K\). Using \(K\leq2+L/512\) gives
\[
28K+\frac{16K^2}{L}\leq122+\frac L{16}
=123+\frac{\log d}{16},
\qquad
(D+1)^4e^{24K+16K^2/L}\leq e^{123}d^{1/16}.
\]
Consequently the conditional squared bounded-variation norm is at most \(4c_{\mathrm{coeff}}^2e^{123}d^{1/16}S_j^2\). Every radius is positive, so \(S_j>0\). Apply \(|\Pi_{[-t,t]}(x)-x|\leq x^2/t\) with \(t=d^{1/4}S_j\). Uniformly in \(\lambda\),
\[
\mathbb E\!\left[|T_j(\lambda)-Z_j(\lambda)|\mid N'\right]
 \leq 4c_{\mathrm{coeff}}^2e^{123}S_jd^{-3/16}
 \qquad\text{on }G_j.
\]
% lean: CausalSmith.Stat.DiscreteBudgetvalueCurve.goodPilot_localJacksonFourPoissonFactorialPath_sq_integral_le The factorial-moment identity also gives
\[
\mathbb E[Z_j(\lambda)\mid N']=P_{j,\lambda}(q_j).
\]
To bound its approximation error, put \(\delta=L/m\). On \(G_j\), the vector \(q_j\) lies in the pilot rectangle, so \cref{obj:lem:jackson-boundary-adaptive} gives
\[
|P_{j,\lambda}(q_j)-f_\lambda(q_j)|
\leq32\{3(1+\epsilon^{-1})+1\}
\left\{
\frac1K\sum_{\zeta\in\mathcal Z}
\sqrt{(q_{\zeta j}-\ell_{\zeta j})(u_{\zeta j}-q_{\zeta j})}
+\frac{S_j}{K^2}\right\}.
\]
The good-pilot inequality and the truncated endpoints imply, coordinatewise,
\[
(q_{\zeta j}-\ell_{\zeta j})(u_{\zeta j}-q_{\zeta j})
\leq8H_0^2q_{\zeta j}\delta.
\]
Indeed, when \(\ell_{\zeta j}>0\), the product is at most \(h_{\zeta j}^2\); the inequalities \(c_{\zeta j}>h_{\zeta j}\) and \(|c_{\zeta j}-q_{\zeta j}|\leq h_{\zeta j}/4\), inserted into the formula for \(h_{\zeta j}\), give the displayed bound. When \(\ell_{\zeta j}=0\), that formula and \(c_{\zeta j}\leq h_{\zeta j}\) give \(u_{\zeta j}\leq8H_0^2\delta\), so the product is at most \(q_{\zeta j}u_{\zeta j}\). Since each \(q_{\zeta j}\leq p_j\) and there are four coordinates,
\[
\sum_{\zeta\in\mathcal Z}
\sqrt{(q_{\zeta j}-\ell_{\zeta j})(u_{\zeta j}-q_{\zeta j})}
\leq4\sqrt{8H_0^2p_jL/m}.
\]
Thus the approximation contribution has the asserted \(\sqrt{p_jL/m}/K+S_j/K^2\) scale, including when \(p_j=0\). Combining it with the conditional clipping contribution gives, on \(G_j\),
% lean: CausalSmith.Stat.DiscreteBudgetvalueCurve.goodPilot_localJacksonPolynomial_boundary_sharp
\[
\sup_{\lambda\in[0,1]}
\left|\mathbb E\{T_j(\lambda)\mid N'\}-f_\lambda(q_j)\right|
 \leq c_\epsilon\left\{
 S_jd^{-3/16}
 +\frac{\sqrt{p_jL/m}}K+\frac{S_j}{K^2}\right\}.
\]
The radius definition and truncation at zero give the pointwise estimates
\[
R_{\zeta j}^{2}
\leq2H_0^2\left\{\frac{N'_{\zeta j}L}{m^2}
+\left(\frac Lm\right)^2\right\},
\qquad
S_j^2\leq4\sum_{\zeta\in\mathcal Z}R_{\zeta j}^2.
\]
Taking expectations and using \(\mathbb E N'_{\zeta j}=m q_{\zeta j}\) and \(\sum_\zeta q_{\zeta j}=p_j\) yields
\[
\mathbb E S_j^2
\leq8H_0^2\left\{\frac{p_jL}{m}
+4\left(\frac Lm\right)^2\right\}
\leq32H_0^2v_j^2.
\]
This is the second-moment bound in \cref{obj:lem:pilot-radius-second-moment}; Cauchy–Schwarz also gives \(\mathbb E S_j\leq\sqrt{32}\,H_0v_j\).
% lean: CausalSmith.Stat.DiscreteBudgetvalueCurve.pilotRadius_sum_sq_integral_le_cellScale Finally,
\(K^{-1}\leq768/L\) and \(L^2d^{-3/16}\leq(35/3)^2\) for \(d\geq16\). After averaging over pilots, these inequalities convert the displayed good-pilot bound to the scale \(a_j\).

For the complementary pilot event, put
\[
B_j=\sum_{\zeta\in\mathcal Z}
 \bigl(|c_{\zeta j}-q_{\zeta j}|+h_{\zeta j}\bigr).
\]
Applying \cref{obj:lem:pilot-bad-score-moment} with \(t=2\) gives the weighted second-moment estimate
\[
\mathbb E[\mathbf1_{G_j^c}B_j^2]
\leq 16\cdot10^{240}e^{-20L}v_j^2
=:c_{\mathrm{mom}}e^{-20L}v_j^2.
\]
% lean: CausalSmith.Stat.DiscreteBudgetvalueCurve.pilotCell_badScore_sq_integral_le
The pilot endpoints give \(|c^0_{\zeta j}-c_{\zeta j}|\leq h_{\zeta j}/2\) and \(R_{\zeta j}\leq h_{\zeta j}\). Clipping bounds \(|T_j(\lambda)-f_\lambda(c_j^0)|\) by \(d^{1/4}S_j\), while \cref{obj:lem:bv-lipschitz} bounds \(\sup_\lambda|f_\lambda(c_j^0)-f_\lambda(q_j)|\) by \(c_{\mathrm{BV}}\|c_j^0-q_j\|_1\). Hence, on \(G_j^c\),
\[
\sup_{\lambda\in[0,1]}|T_j(\lambda)-f_\lambda(q_j)|
 \leq(c_{\mathrm{BV}}+d^{1/4})
 \sum_{\zeta\in\mathcal Z}
 \bigl(|c_{\zeta j}-q_{\zeta j}|+h_{\zeta j}\bigr).
\]
Since \(d^{1/4}\geq1\), taking square roots yields, for a positive \(c_{\mathrm{env}}\),
\[
\left\|\mathbf1_{G_j^c}
 \sup_{\lambda\in[0,1]}|T_j(\lambda)-f_\lambda(q_j)|
 \right\|_{L_2}
 \leq c_{\mathrm{env}}d^{1/4}e^{-10L}v_j.
\]
Splitting the expectation over \(G_j\) and \(G_j^c\), and using the preceding conditional and bad-event bounds, gives a positive \(c_{\mathrm{bias}}\) such that
\[
\sup_{\lambda\in[0,1]}
 \left|\mathbb ET_j(\lambda)-f_\lambda(q_j)\right|
 \leq c_{\mathrm{bias}}a_j.
\]
In particular, the exponentially weighted bad-event contribution is absorbed at this scale because \(v_j\leq L^2a_j\) and
\(d^{1/4}e^{-10L}L^2\) is bounded for \(d\geq16\). % lean: CausalSmith.Stat.DiscreteBudgetvalueCurve.uncappedJacksonCellBias_le_rate
% lean: CausalSmith.Stat.DiscreteBudgetvalueCurve.badPilotCellEnvelope_le

Define the good- and bad-pilot error paths by
\[
W_j^{\mathrm{good}}(\lambda)
 =\mathbf1_{G_j}\{T_j(\lambda)-f_\lambda(q_j)\},
\qquad
W_j^{\mathrm{bad}}(\lambda)
 =\mathbf1_{G_j^c}\{T_j(\lambda)-f_\lambda(q_j)\}.
\]
The pilot and evaluation counts are independent across cells. On \(G_j\), clipping is one-Lipschitz and cannot increase total variation. The preceding conditional factorial-moment bound, the coefficient envelope, and \cref{obj:lem:bv-lipschitz} therefore give
\[
\mathbb E\!\left[\|W_j^{\mathrm{good}}\|_{BV([0,1])}^2\mid N'\right]
 \leq c\,\mathbf1_{G_j}d^{1/16}S_j^2.
\]
Indeed, the clipped factorial path has squared bounded-variation norm at most a constant times \(d^{1/16}S_j^2\); its correction from the pilot center to \(q_j\) has bounded-variation norm at most \(c_\epsilon S_j\), since \(q_j\) lies in the pilot rectangle. The second-moment calculation above gives \(\mathbb ES_j^2\leq32H_0^2v_j^2\). Consequently,
\[
\mathbb E\|W_j^{\mathrm{good}}\|_{BV([0,1])}^2
 \leq c\,d^{1/16}v_j^2.
\]
The paths are independent across cells, continuous, and have the finite variation and second moments established above. By \cref{obj:lem:bv-maximal-cell-paths}, and because \(\|W_j^{\mathrm{good}}\|_\infty+\operatorname{TV}(W_j^{\mathrm{good}})\leq2\|W_j^{\mathrm{good}}\|_{BV([0,1])}\),
\[
\begin{aligned}
\mathbb E\left\|\sum_j
 (W_j^{\mathrm{good}}-\mathbb EW_j^{\mathrm{good}})
 \right\|_\infty^2
&\leq 16384\sum_j\mathbb E
 \left[\bigl(\|W_j^{\mathrm{good}}\|_\infty+
 \operatorname{TV}(W_j^{\mathrm{good}})\bigr)^2\right]\\
&\leq65536\sum_j\mathbb E
 \|W_j^{\mathrm{good}}\|_{BV([0,1])}^2\\
&\leq c\,d^{1/16}
 \left(\frac{2L}{m}+2d\frac{L^2}{m^2}\right).
\end{aligned}
\]
The last scale follows from \(\sum_jp_j=1\) and
\(\sum_jv_j^2\leq2L/m+2d(L/m)^2\). On the branch \(d/L\leq n=8m\), the inequalities \(L\leq17d^{1/16}\) and \(d/m\leq8L\) turn the last display into
\[
\mathbb E\left\|\sum_j
 (W_j^{\mathrm{good}}-\mathbb EW_j^{\mathrm{good}})
 \right\|_\infty^2
 \leq c_{\mathrm{good}}\frac d{mL}.
\]
% lean: CausalSmith.Stat.DiscreteBudgetvalueCurve.centeredIdealGoodPilotErrorRisk_le_rate

For the bad-pilot paths, centering, the triangle inequality, and Cauchy–Schwarz give
\[
\mathbb E\left\|\sum_j
 (W_j^{\mathrm{bad}}-\mathbb EW_j^{\mathrm{bad}})
 \right\|_\infty^2
 \leq4d\sum_j\mathbb E\|W_j^{\mathrm{bad}}\|_\infty^2
 \leq c\,d^{3/2}e^{-20L}
 \left(\frac{2L}{m}+2d\frac{L^2}{m^2}\right).
\]
Here \(L\leq d\), \(dL/m\leq8L^2\), and
\(d^{1/2}e^{-20L}\leq d^{-4}\). They absorb the remaining factors and give a positive \(c_{\mathrm{bad}}\) with
\[
\mathbb E\left\|\sum_j
 (W_j^{\mathrm{bad}}-\mathbb EW_j^{\mathrm{bad}})
 \right\|_\infty^2
 \leq c_{\mathrm{bad}}\frac d{mL}.
\]
% lean: CausalSmith.Stat.DiscreteBudgetvalueCurve.centeredIdealBadPilotErrorRisk_le_rate

It remains to combine the centered terms with their mean. Define
\[
\mathfrak b=c_{\mathrm{bias}}\sum_j a_j.
\]
The good and bad error paths partition \(\sum_j(T_j-f_\cdot(q_j))\), so the absolute value of its mean at every \(\lambda\) is at most \(\mathfrak b\). Since \(\sum_jp_j=1\), Cauchy–Schwarz and \(d/(nL)\leq1\) give
\[
\mathfrak b^2
 \leq144c_{\mathrm{bias}}^2\frac d{nL}.
\]
Applying \((x+y+z)^2\leq3(x^2+y^2+z^2)\) to the centered good path, centered bad path, and mean, then using \(m=n/8\), proves
\[
\mathbb E_{\mathrm{id}}
 \left\|\sum_jT_j-F_P\right\|_\infty^2
 \leq\Gamma_{\mathrm{id}}\frac d{nL},
\qquad
\Gamma_{\mathrm{id}}
 =24c_{\mathrm{good}}+24c_{\mathrm{bad}}
   +432c_{\mathrm{bias}}^2.
\]
% lean: CausalSmith.Stat.DiscreteBudgetvalueCurve.idealProcessRisk_le_rate

For the averaged process of \cref{obj:def:jf-estimator}, define the nonnegative count-table function
\[
\mathcal G(N',N)
 =\left\|\sum_jT_j(\cdot;N',N)-F_P\right\|_\infty^2.
\]
Its ideal expectation is finite by the preceding process bound. Conditional Jensen bounds the averaged-process risk by the capped auxiliary risk. Apply \cref{obj:lem:capped-marked-count-identity} to \(\mathcal G\); in the uncapped marked-Poisson experiment, \(M\) is the total number of marked records, and the count-table marginal is the ideal law. Hence
\[
\mathbb E_{O,\mathrm{aux}}[\mathbf1_{\{M\leq n\}}\mathcal G(N',N)]
 =\mathbb E_{\mathrm{uncap}}[\mathbf1_{\{M\leq n\}}\mathcal G(N',N)]
 \leq\mathbb E_{\mathrm{id}}\mathcal G(N',N).
\]
On \(M>n\), the auxiliary process is zero; \(0\leq f_\lambda(q_j)\leq p_j\) implies \(\|F_P\|_\infty\leq1\). The Poisson cap-tail bound therefore gives
\[
\begin{aligned}
\mathbb E\|\widehat F-F_P\|_\infty^2
&\leq\mathbb E_{O,\mathrm{aux}}\|\widetilde F^{\mathrm{cap}}-F_P\|_\infty^2\\
&=\mathbb E_{O,\mathrm{aux}}[\mathbf1_{\{M\leq n\}}\mathcal G(N',N)]
  +\Pr(M>n)\|F_P\|_\infty^2\\
&\leq\mathbb E_{\mathrm{id}}
       \left\|\sum_jT_j-F_P\right\|_\infty^2
       +e^{-n/16}.
\end{aligned}
\]
The elementary bounds \(e^{-n/16}\leq16/n\) and
\(L=1+\log d\leq d\) show that \(e^{-n/16}\leq16d/(nL)\). Hence
\[
\mathbb E\|\widehat F-F_P\|_\infty^2
 \leq(\Gamma_{\mathrm{id}}+16)\frac d{nL}.
\]
Under \cref{obj:ass:iid-sampling}, the product law used in this expectation is the law of the observed sample. % lean: CausalSmith.Stat.DiscreteBudgetvalueCurve.capTransferRisk
% lean: CausalSmith.Stat.DiscreteBudgetvalueCurve.exponential_cap_remainder_le_rate

Choose one constant larger than every coefficient used above; for example, with the cap-transfer coefficient written \(c_{\mathrm{cap}}=1\), take
\[
C_\epsilon
 =1+c_{\mathrm{coeff}}+e^{115}
   +8c_{\mathrm{bias}}+8c_{\mathrm{good}}
   +c_{\mathrm{env}}+c_{\mathrm{bad}}
   +\Gamma_{\mathrm{id}}
   +c_{\mathrm{cap}}(\Gamma_{\mathrm{id}}+16).
\]
All terms are nonnegative, so \(C_\epsilon>0\). The five estimates established for arbitrary observed categorical laws and every pilot realization give \(C_\epsilon\in\mathsf H_{n,d}\) as defined in \cref{obj:def:process-handle}.

Finally, take the observed margin of the full-data law \(P\in\mathcal M_{d,\epsilon}\) from \cref{obj:def:model-class}. Summing its four observed category masses in cell \(j\) gives the same covariate mass \(p_j\) as summing the corresponding full-data masses. Moreover,
\[
\sqrt{\frac{p_j}{mL}}+\frac1{mL}
 =\sqrt{\frac{8p_j}{nL}}+\frac8{nL}
 \leq8\left\{\sqrt{\frac{p_j}{nL}}+\frac1{nL}\right\},
\qquad
\frac d{mL}=8\frac d{nL}.
\]
Thus the common choice \(C_\epsilon\) gives the stated \(nL\)-scale cell-bias and good-pilot bounds, the stated \(mL\)-scale bad-pilot bound, the cellwise bad-pilot envelope, and both process-risk bounds simultaneously. % lean: CausalSmith.Stat.DiscreteBudgetvalueCurve.integrated_process_certificate
\end{proof}
